problem:  
Let \( (X,\omega_0) \) be a compact complex threefold endowed with a balanced Hermitian metric (\(d(\omega_0^2)=0\)) and let \(E \to X\) be a stable holomorphic bundle equipped with a Hermitian–Yang–Mills metric \(h\).  
Consider conformal deformations \(\omega = e^{u}\omega_0\) and the **Strominger anomaly equation:**  
\[
i\partial\bar\partial \omega = \alpha'(\operatorname{Tr}(R^B\wedge R^B) - \operatorname{Tr}(F_h\wedge F_h)),
\]
where \(R^B\) is the curvature of the Bismut connection of \((X,\omega)\) and \(F_h\) is the curvature of the Chern connection on \(E\).  
Assume the required cohomological condition \(c_2(TX)=c_2(E)\in H^4(X,\mathbb{R})\).  

**Problem:**  
Known existence theorems (Fu–Yau, Phong–Sturm, Fei–Picard–Zhang) produce solutions primarily for **\(\alpha'>0\)** under special geometric assumptions.  
Is the **sign of \(\alpha'\)** intrinsically constrained by the torsion class of the balanced metric, or can there exist a compact *balanced but non–Kähler* \(X\), not conformally Calabi–Yau, together with a bundle \(E\to X\) satisfying \(c_2(TX)=c_2(E)\), such that smooth conformal factors \(u_+\) and \(u_-\) produce solutions of the anomaly equation for **both \(\alpha'>0\)** and **\(\alpha'<0\)** within the same conformal class?  
If such dual solutions exist, does the sign reversal correspond to a geometric operation on the Bismut torsion \(T=Jd\omega\) (for example, inversion of the Lee form or of the (2,1)+(1,2) component of torsion)?  

Why is it a "good" problem:  
This version is now consistent with the standard formulation of the Strominger system and incorporates the necessary topological constraint. It focuses on the **sign dependence of the anomaly coupling \(\alpha'\)**, a structural aspect that remains analytically and geometrically unresolved.  
Existing literature shows solvability for specific sign regimes (e.g., \(\alpha'>0\), with ellipticity depending on the nonnegativity of certain curvature operators), but no general principle is known linking sign(\(\alpha'\)) to the torsion geometry of balanced manifolds.  

Addressing this question requires combining:  
- **Analytic insight:** understanding how sign(\(\alpha'\)) modifies the elliptic structure of the Fu–Yau–type PDE and torsion-curvature coupling;  
- **Geometric insight:** interpreting the torsion (Lee form, intrinsic SU(3)–structure) as a bridge between positive and negative flux regimes;  
- **Topological coherence:** ensuring anomaly cancellation via \(c_2(TX)=c_2(E)\).  

If the same background admits solutions for both signs, this would indicate an unexpected **torsion duality** within non–Kähler Hermitian geometry; if not, it would reveal a **topological or analytic obstruction** that classifies flux compatibilities in terms of torsion data.  

The statement achieves simplicity of form while embedding deep analytic and geometric questions, crossing complex geometry, gauge theory, and nonlinear PDE analysis—meeting the “narrow entrance, profound depth” standard for a good research-level problem.